\documentclass[11pt,a4paper]{article}
\usepackage[T1]{fontenc}
\usepackage{lmodern,amsmath,amssymb,microtype,booktabs}
\usepackage[margin=25mm]{geometry}
\usepackage[colorlinks=true,urlcolor=blue,linkcolor=blue]{hyperref}
\newcommand{\Rh}{\widehat R}
\newcommand{\dd}{\,\mathrm d}
\title{Physical spectral weights in the local dS$_3$ crossing limit}
\author{Alan Day \and Chengheng\thanks{AI research assistant (OpenAI).}}
\date{Working research note --- 31 October 2026}
\begin{document}\maketitle
\begin{abstract}
We insert the actual regulated creation amplitudes into the previously
derived local scalar crossing limit. The representation-label change of
variables has unit absolute Jacobian. After the physical normalization
is included exactly once, the weights form a uniformly bounded $C^2$
family on a fixed compact transverse window. This proves a local,
smoothly cut off physical pairing limit. It does not remove the cutoff
or prove the full Bose free-observer limit. Numerical checks reuse the
saved complete-kernel nodes and independently assemble their weights
from the original vertices and four-body norms.
\end{abstract}

\section{State, measure, and exact weight}
Keep $\ell=1$, $G_N=0$, proper $SO^+(3,1)$, two distinguishable BD
scalar observers with labels at least a fixed $s>0$, KG matter
$p,q>0$, unit local coupling, and the half-period regulated creation
blocks. Time is generated by $h(x)=\sqrt{1+x^2}$.
The initial preparation here is conditional label $c=M$, defined as
the limit of normalized narrow mass packets. It is not a fixed initial
packet sent to large mass without changing its preparation.

Our earlier channel conventions \cite{two,packet} give the coefficient
of $U_{00}(b,d;a,q,c,p)$ in the crossed pairing as
\begin{align}
 A_M&=\frac{g(a)}{g(c)}
 \sqrt{\Rh_q(a,b)\Rh_p(b,c)\Rh_p(a,d)\Rh_q(d,c)},\label{amp}\\
 \Phi_M(t)&=\exp\{it[h(a)+h(c)-h(b)-h(d)]\},\quad F_M=A_M\Phi_M,\label{phase}\\
 g(x)&=\frac{x e^{\pi[h(x)-x]}}{1+e^{-2\pi x}},\qquad
 \Rh_p(x,y)=\frac{A-B}{8\pi(A+C)(B+C)},\label{hat}\\
 &A=\cosh\pi(x+y),\quad B=\cosh\pi(x-y),\quad C=\cosh\pi p.\nonumber
\end{align}
In these conventions the conditional integration measure is
$\dd a\dd b\dd d$. Before conditioning it is
$\pi_0(c)\dd c\dd a\dd b\dd d$, where
$\pi_0(c)=\rho(c)|f(c)|^2$, $\rho(c)=2c^2/\pi$.

The same amplitude, before cancellation, is
\[
 A_M=\frac{\overline{E_{abq}E_{bcp}}E_{adp}E_{dcq}
             e^{\pi[h(a)-h(c)]}}
 {\rho(c)\sqrt{W_sW_t}},\qquad
 W_s=\rho(b)K_{abq}K_{bcp},\quad
 W_t=\rho(d)K_{adp}K_{dcq}.
\]
The fixed spherical vertices are real positive in this convention.
These equations already contain the observer and CG normalizations;
multiplying another $\rho(b)\rho(d)$ into the integral would be incorrect.

\newpage
\section{Coordinates and a uniform weight lemma}
Set
\[
 c=M,\quad b=M+u,\quad a=M+u+v,\quad d=M+v+r.
\]
Both $|\partial(a,b,c,d)/\partial(M,u,v,r)|$ and the conditional
$|\partial(a,b,d)/\partial(u,v,r)|$ equal one.
If one instead uses energies $\epsilon_x=h(x)$, then
\[
 \dd a\dd b\dd d=
 \prod_{x=a,b,d}\frac{h(x)}{x}\,
 \dd\epsilon_a\dd\epsilon_b\dd\epsilon_d .
\]
The factor tends to one with $O(M^{-2})$ error on compact transfers,
but it is not identically one. Conditioning in energy uses the
transformed probability density $\pi_0(c)h(c)/c$.

Define $k_p(x)=[8\pi(\cosh\pi x+\cosh\pi p)]^{-1}$ and
\begin{equation}
 F_\infty(u,v,r;t)=
 \sqrt{k_q(v)k_p(u)k_p(u-r)k_q(v+r)}\,e^{-itr}.
 \label{limitweight}
\end{equation}
On any fixed compact set of real $(u,v,r)$ and positive $(p,q)$,
uniformly for $|t|\le T$, the exact weights satisfy
\begin{equation}
 F_M=\left(1+\frac{u+v}{M}\right)F_\infty
       +O_{C^2_r,T}(M^{-2}).                         \label{C2}
\end{equation}
All masses exceed $s$ for sufficiently large $M$ on this set.
Here $C^2_r$ controls the function and its first two $r$ derivatives;
it is not an assertion for growing support or growing time.

\paragraph{Proof.}
The exact identity
\[
 \Rh_p(x,y)=k_p(x-y)
 \left[1-\frac{\cosh\pi(x-y)+\cosh\pi p}
                  {\cosh\pi(x+y)+\cosh\pi p}\right]
\]
has bracket $1+O_{C^2_r}(e^{-2\pi M})$ for $x=M+O(1)$,
$y=M+O(1)$. Derivatives of its numerator are bounded on the compact
set; derivatives of the reciprocal denominator retain the exponential
bound. The bracket is bounded away from zero for large $M$, so taking
the positive square root preserves the estimate.

Write $\delta(x)=h(x)-x=1/(h(x)+x)$. Its derivatives obey
\[
 |\delta'(x)|\le\frac1{2x^2},\qquad
 \delta''(x)=(1+x^2)^{-3/2}\le x^{-3}\quad(x>0).
\]
Consequently
$h(a)+h(c)-h(b)-h(d)=-r+O_{C^2_r}(M^{-2})$.
Also
\[
 \frac{g(M+u+v)}{g(M)}
 =\left(1+\frac{u+v}{M}\right)
 e^{\pi[\delta(M+u+v)-\delta(M)]}
 \frac{1+e^{-2\pi M}}{1+e^{-2\pi(M+u+v)}}
 =1+\frac{u+v}{M}+O(M^{-2}).
\]
This factor is independent of $r$; the derivative bounds for the phase,
the product rule, and the positivity of the compact limiting weights
prove \eqref{C2}. No differentiation of an unspecified remainder is used.

\newpage
\section{A local physical pairing theorem}
Let $R>0$ be fixed and let $\chi(u,v,r)$ be a fixed $C^2$ compactly
supported cutoff, supported in $r\in[-R,R]$ and a compact transfer set
on which
\begin{equation}
 \max_j|\alpha_j|+R/2\le1/2,\qquad
 \alpha=\tfrac12(v+q,p-u,-u-p,v-q).                 \label{window}
\end{equation}
This sufficient small-window restriction comes from the earlier kernel
estimate; the weight lemma itself does not require it.
Define the compact part of the actual conditional pairing,
\[
 I_M^\chi(t)=\int\chi(u,v,r)F_M(u,v,r;t)
                   U_{00}(M;u,v,r)\dd u\dd v\dd r .
\]
The preceding local-kernel note \cite{local} gives, uniformly on this set,
\[
 U_{00}=\frac{\sin(r\log M-\vartheta(u,v,r))}{\pi r}
                +O(M^{-1}),\qquad \vartheta(u,v,0)=0 ,
\]
including the removable value at $r=0$. Its Dirichlet estimate applies
uniformly to fixed-support test families with bounded $C^2_r$ norms.
The lemma supplies those bounds for $\chi F_M$, and hence
\begin{equation}
 I_M^\chi(t)=
 \int \chi(u,v,0)k_p(u)k_q(v)\dd u\dd v
 +O_{\chi,T}(1/\log M)+O_{\chi,T}(1/M).            \label{physical}
\end{equation}
The convergence is uniform on compact real times.
To see the last integration explicitly, apply the one-dimensional
estimate for each $(u,v)$; the cutoff's support has finite area and the
constants are uniform there. At $r=0$, \eqref{C2} gives
$F_M=k_p(u)k_q(v)+O(M^{-1})$ uniformly. This completes the local argument.

This result treats a compact contribution to the original observable.
The cutoff is an analysis device; we have not replaced the original
preparation or internal spectrum by a filtered physical model.
The leading time independence applies to the crossed phase
\eqref{phase}. It does not say every Bose exchange matrix element has
a time-independent limit.

For equal species, the exchange phase
$\exp\{it[h(d)-h(b)]\}$ instead tends to $e^{it(v+r-u)}$.
The same local argument then gives the integrand
$\chi(u,v,0)k_p(u)k_p(v)e^{it(v-u)}$.
If a future global theorem permits removal of all cutoffs, these two
expressions would give $L_p(0)^2$ and $|L_p(t)|^2$, respectively, with
\[
 L_p(t)=\frac{\sin(pt)}{4\pi\sinh(\pi p)\sinh t},
 \qquad L_p(0)=\frac{p}{4\pi\sinh(\pi p)}.
\]
They are the two different terms needed in the Bose Gram permanent.
At present, replacing $\chi$ by one is unproved: the restricted domain
\eqref{window} cannot exhaust the full transfer plane, and the original
oscillatory pairing need not be absolutely integrable. A global argument
must combine a larger compact-kernel theorem with the physical Hilbert
norm tail bounds. In particular $U_{00}$ being a contraction matrix
kernel does not imply the pointwise bound $|U_{00}|\le1$.

\newpage
\section{Numerical evidence and continuation}
The new diagnostics use 85-digit arithmetic and preserve full complex
values. Twelve points at $M=4,16,64,256$ compare \eqref{amp} with the
independent $E/W/\rho$ expression. The largest relative difference is
about $2.95\,10^{-83}$, against the recorded empirical tolerance
$10^{-65}$. The fixed samples include unequal species and nonzero
transfers; they do not prove uniformity.

Seventy-two derivative comparisons check analytic derivatives of
\eqref{limitweight} against numerical differentiation, for both crossed
and exchange phases and orders $0,1,2$. The corresponding finite-$M$
derivatives and $M^2$-scaled remainders in \eqref{C2} are saved as
observations, without imposing a convergence threshold. Their largest
moduli at the four masses are approximately
$0.000474$, $0.000506$, $0.000508$, $0.000508$.

For the physical local integral at $(u,v,p,q)=(0.1,-0.1,0.4,0.4)$,
use $\chi_r=(1-(r/R)^2)^3$ on $[-R,R]$, zero outside, $R=0.25$.
The following are ratios to the limiting value $k_p(u)k_q(v)$, using
the 48-node rule:
\begin{center}
\begin{tabular}{rrr}
\toprule
$t$ & $M=64$ & $M=256$\\
\midrule
0 & $0.470982$ & $0.552526$\\
1 & $0.469502-0.001076i$ & $0.550854-0.001179i$\\
2 & $0.465087-0.002135i$ & $0.545862-0.002340i$\\
\bottomrule
\end{tabular}
\end{center}
These values are far from the limiting ratio one. In particular the
small imaginary component need not decrease monotonically at these
masses. The 24/48 rules reuse all 144 complete complex kernel nodes
from \cite{local}; they are not 144 new evaluations of $U_{00}$.
Twelve reassemblies of the weighted integrals agree with the independent
vertex representation. One exact symbolic Jacobian check completes the
97 checks. All propagated series tails exclude quadrature and rounding.

The implementation is \texttt{code/ds3\_physical\_local.py}; diagnostics
and parameters are in
\texttt{results/stage4/ds3-physical-local/diagnostics.json}.
The driver refuses to overwrite existing results. Input/output hashes,
dependency versions and the compilation review accompany this note.
The proof is the compact estimate above, not the finite sample count.
The next research task is to extend the compact kernel estimate beyond
\eqref{window}, then control the complement in physical Hilbert norm.
A fixed initial packet, a moving threshold, and growing observation
time each require separate arguments. Full Fock domains, finite-$G_N$
dynamics and a complete microscopic holographic dual remain open.

\begin{thebibliography}{9}
\bibitem{two} Alan Day and Chengheng,
\emph{dS$_3$ two-matter sector and noncrossed limit},
project note \texttt{ds3-two-matter}, 31 October 2026.
\bibitem{packet} Alan Day and Chengheng,
\emph{Physical crossed packets}, project note \texttt{ds3-crossed-packet},
31 October 2026.
\bibitem{local} Alan Day and Chengheng,
\emph{A uniform transverse limit of the dS$_3$ crossing kernel and its
local action on test functions}, project note \texttt{ds3-uniform-local},
31 October 2026.
\bibitem{lprs} J. Liu, E. Perlmutter, V. Rosenhaus and D. Simmons-Duffin,
\href{https://arxiv.org/html/1808.00612v3}{arXiv:1808.00612v3},
Eqs.\ (3.35)--(3.38), underlying crossing kernel.
\bibitem{wilson} W. Groenevelt,
\href{https://arxiv.org/html/math/0306424v1}{arXiv:math/0306424v1},
Proposition 4.7, precedent for the removable-pole Dirichlet argument.
\end{thebibliography}
\end{document}